\documentclass{article}

% [preprint] = non-anonymized, for circulation outside review
\usepackage[preprint]{neurips_2026}

\usepackage[utf8]{inputenc}
\usepackage[T1]{fontenc}
\usepackage{hyperref}
\usepackage{url}
\usepackage{booktabs}
\usepackage{amsmath}
\usepackage{microtype}

\title{Which Verifier Catches What? Mapping Error Severity to Verifier Strategy in Multi-Agent LLM Pipelines}

\author{%
  Tilak Parajuli\\
  SISTER 2026, AI/ML Track\\
  Dubai Computer Science Society \& Delta Rising Foundation\\
  Kathmandu, Nepal\\
  \texttt{tilak.parajuli.58@gmail.com}
   \And
  Johnny Kozman\\
  SISTER 2026, AI/ML Track\\
  Dubai Computer Science Society \& Delta Rising Foundation\\
  Cairo, Egypt\\
  \texttt{kozmanjohnny82@gmail.com}
   \And
  Arushi Waddepalli\\
  SISTER 2026, AI/ML Track\\
  Dubai Computer Science Society \& Delta Rising Foundation\\
  Hyderabad, India\\
  \texttt{arushiwaddepalli09@gmail.com}
}

\begin{document}

\maketitle

\paragraph{Area.} Research: multi-agent LLM reliability and error propagation. Method: controlled fault injection across severity levels and verifier prompt strategies (direct, CoT, rubric).


\begin{abstract}
Multi-agent LLM pipelines (plan, solve, verify) are rapidly being deployed,
yet how errors travel through them is not quantitatively understood. Existing
work taxonomizes failures observed in the wild \citep{cemri2025mast} but does
not measure propagation causally. We adapt fault injection, a standard method
from software reliability engineering, to LLM pipelines: we corrupt one
intermediate output with a known, scripted perturbation (a number swap or an
operation flip), graded into three severity levels (mild, moderate, severe),
and measure whether the fault is caught, corrected, or silently propagated to
the final answer. On a fixed sample of roughly 200 GSM8K problems, extending
to HotpotQA, we compare six conditions: a single-agent baseline, a clean
three-agent pipeline, early- and late-stage injection, and verifier-removed
ablations of both. Because every stage output is stored, we additionally
replay the verification stage over identical traces under three prompt
strategies (direct, chain-of-thought, rubric checklist), so that differences
in catch rate are attributable to the strategy alone; the result is a map of
which verifier strategy catches which error type at which severity. Outcomes
are analyzed with paired statistics: McNemar's exact test, bootstrap
confidence intervals, Wilson intervals, and logistic regression on fault
survival. The study is fully executable in four weeks on API access; all
code, prompts, and traces will be released publicly.
\end{abstract}

\section{Introduction}

The dominant pattern in applied LLM systems is no longer a single model call
but a pipeline of specialized agents: one decomposes the task, another
executes it, a third verifies the result. The implicit promise of this design
is robustness; later stages are expected to absorb the mistakes of earlier
ones. The literature is split on whether that promise holds. Multi-agent
debate has been reported to improve factuality \citep{du2023debate}, while
self-correction studies find that models largely cannot identify their own
reasoning errors without external feedback \citep{huang2023selfcorrect}. The
most closely related work, the MAST taxonomy \citep{cemri2025mast},
catalogues fourteen failure modes observed across seven deployed multi-agent
frameworks. As an observational study, however, it cannot say how often an
early-stage error survives to the output, nor what a verification stage
causally contributes. Meanwhile, mitigation layers that gate or govern
inter-agent messages are already being engineered on the assumption that such
errors cascade \citep{yoffe2024debunc, xie2026spark}; measuring that
assumption is the missing step.

We close that gap experimentally. Because each fault is planted by us, ground
truth about the origin of every downstream failure is known by construction,
turning ``agents sometimes fail'' into a measurable, attributable quantity.
This study aims to answer four specific research questions:
\begin{itemize}
    \item \textbf{RQ1 (Error Propagation \& Position):} How often do injected
    faults propagate to the final output, and does this survival rate differ when
    errors are introduced in the planning stage (Decomposer) versus the execution
    stage (Solver)?
    \item \textbf{RQ2 (Severity \& Sensitivity):} How does error severity
    (mild, moderate, severe) affect downstream propagation? Do verifiers only
    catch severe, fatal errors, or are they sensitive to subtle numerical slips?
    \item \textbf{RQ3 (Prompting Strategies \& Efficiency):} How do Direct,
    Chain-of-Thought (CoT), and Rubric prompting strategies compare in catch rate,
    verifier correction precision, and propagation control? What is the token cost
    overhead associated with more structured prompting? Does the relative advantage
    of each prompting strategy vary depending on fault type (number swap, operation
    flip, semantic logic error) and severity (mild, moderate, severe)?
    \item \textbf{RQ4 (Causal Map \& Error Types):} How do verifier presence
    (ablation) and prompt strategy interact with specific error types (number swaps,
    operation flips, and semantic logic errors)? Which verifier configuration is
    best suited for catching each type of reasoning failure?
\end{itemize}

The design meets the program's stated criteria directly: public data with
defined ground truth (GSM8K, HotpotQA); a simple baseline (single agent);
metrics fixed before experimentation (accuracy, propagation rate, catch and
correction rate, cost per solved problem); a scope achievable in two to four
weeks; novelty in the evaluation method rather than a demo; and evidence that
includes error analysis and stated limitations.

\bibliographystyle{plainnat}
\bibliography{references}

\end{document}
